\section{Blackwell Ownership and Memory Architecture}

\subsection{The copied V382 schedule is the control}

The low-precision tree is a namespace-isolated copy of the optimized
ThunderKittens BF16 V382 backward \cite{thunderkittens}.  The BF16 source is
kept as the invariant control: low precision changes representations and
selected producer/consumer handoffs inside the same causal scheduler rather
than replacing it with a small standalone matrix benchmark.

The retained specialization uses one two-CTA cluster.  Each CTA has 512
threads divided into loader, tensor-issue, statistics, compute, and reducer
roles.  A K/V owner keeps its $dK$ and $dV$ accumulators resident while it
visits the legal query tiles.  The two CTAs exchange the fragments needed to
form full dQ and dK operands.  This owner topology is the reason $dK$ and
$dV$ are single-writer outputs while dQ is a cross-owner reduction.

\begin{table}[htbp]
\centering
\caption{Inherited schedule versus the active low-precision changes.}
\label{tab:inherited-vs-lowp}
\small
\begin{tabularx}{\textwidth}{p{.23\textwidth}X X}
\toprule
Component & Inherited from BF16 V382 & Active low-precision change \\
\midrule
Ownership & two-CTA K/V owner; persistent causal query loop &
unchanged \\
Tile geometry & M128$\times$N128, $D_{QK}=192$, $D_V=128$ &
unchanged \\
Score & three K64 reductions into FP32 TMEM &
NVFP4 Q/K operands and one adaptive scale page \\
P and dS & recompute P from saved LSE; produce both dS orientations &
P is E4M3 and reused; dS is emitted directly as E4M3 \\
$dP,dV$ & BF16 tensor operands, FP32 accumulation &
E4M3 tensor operands with exact power-of-two corrections \\
$dQ,dK$ & BF16 tensor operands, FP32 accumulation &
mixed E4M3$\times$E2M1 tensor instructions \\
Publication & cross-owner dQ reduction; single stores for dK/dV &
dQ reduction payload becomes BF16; dK/dV ownership unchanged \\
\bottomrule
\end{tabularx}
\end{table}

\subsection{TMEM is full and heavily aliased}

The active kernel uses all 512 logical TMEM columns.  Long-lived dK and dV
occupy columns 0--319.  Score, P, dP, dS, scale pages, and dQ time-share
columns 320--511.  Figure~\ref{fig:bwd-tmem-map} shows logical bases and the
dominant lifetimes; the lower rows are temporal overlays, not simultaneous
independent allocations.

\begin{figure}[htbp]
\centering
\begin{tikzpicture}[x=0.0195cm,y=.70cm,font=\scriptsize]
  \draw[thick] (0,0) rectangle (512,1);
  \fill[gradientpurple!24] (0,0) rectangle (192,1);
  \fill[outputgreen!28] (192,0) rectangle (320,1);
  \fill[lightgray] (320,0) rectangle (512,1);
  \draw (192,0)--(192,1);
  \draw (320,0)--(320,1);
  \node at (96,.53) {persistent $dK$ [0,192)};
  \node at (256,.53) {persistent $dV$ [192,320)};
  \node at (416,.53) {recycled workspace [320,512)};

  \draw[thick] (320,-1.35) rectangle (512,-.35);
  \fill[scoreblue!30] (320,-1.35) rectangle (448,-.35);
  \fill[scaleyellow!42] (448,-1.35) rectangle (464,-.35);
  \fill[lightgray] (464,-1.35) rectangle (512,-.35);
  \draw (448,-1.35)--(448,-.35);
  \draw (464,-1.35)--(464,-.35);
  \node at (384,-.82) {score [320,448)};
  \node[rotate=90] at (456,-.82) {scale};
  \node at (488,-.82) {scratch};

  \draw[thick] (320,-2.70) rectangle (512,-1.70);
  \fill[probred!28] (320,-2.70) rectangle (352,-1.70);
  \fill[gradientpurple!20] (384,-2.70) rectangle (512,-1.70);
  \draw (352,-2.70)--(352,-1.70);
  \draw (384,-2.70)--(384,-1.70);
  \node at (336,-2.17) {P};
  \node at (368,-2.17) {meta};
  \node at (448,-2.17) {$dP/dS$ base 384};

  \draw[thick] (320,-4.05) rectangle (512,-3.05);
  \fill[lightgray] (320,-4.05) rectangle (416,-3.05);
  \fill[gradientpurple!38] (416,-4.05) rectangle (512,-3.05);
  \draw (416,-4.05)--(416,-3.05);
  \node at (368,-3.52) {released prefix};
  \node at (464,-3.52) {$dQ$ [416,512)};

  \foreach \x in {0,192,320,384,416,448,464,512}
    \node[below] at (\x,-4.12) {\x};
  \node[anchor=east] at (315,.5) {persistent};
  \node[anchor=east] at (315,-.85) {score phase};
  \node[anchor=east] at (315,-2.20) {P/$dP$ phase};
  \node[anchor=east] at (315,-3.55) {$dQ$ phase};
\end{tikzpicture}
\caption{Retained logical TMEM map.  dK spans a 128-column main tile and a
64-column tail; dV spans 128 columns.  The score uses [320,448), its scale page
starts at 448, and the compact dQ accumulator uses [416,512).  P, dP, and dS
reuse retired score/dQ addresses according to phase.}
\label{fig:bwd-tmem-map}
\end{figure}

The most important consequence is negative: another complete M128$\times$N128
FP32 score bank needs 128 columns, but no such bank exists.  Splitting a score
or dQ issue can expose an earlier prefix, yet barriers cannot manufacture
storage.  Experiments that drain half of persistent dV to buy columns lose
more to scratch traffic than they gain from score lookahead.

\subsection{The alias boundary defines legal issue order}

The transient region follows the ownership chain
\begin{equation}
  \text{score}\rightarrow
  \text{P publication}\rightarrow
  dP/dV\ \text{consume}\rightarrow
  dS\rightarrow dQ\ \text{consume}\rightarrow
  \text{next score}.
  \label{eq:bwd-alias-chain}
\end{equation}
The next score cannot overwrite columns 416--447 until the current dQ
consumer has captured that prefix.  Conversely, delaying score issue until
all dQ publication completes exposes unnecessary latency.  The retained
schedule uses issue markers: score and dQ consumers may enqueue their TMEM
loads after the producer instruction has been issued, then let the native
TMEM fence cover completion.  Full completion barriers remain for final
drain and lifetime proof.

\subsection{Shared memory and register boundary}

The profiled hot specialization uses 231,484 bytes of dynamic shared memory,
128 registers per thread, and no local stack or spills.  Shared memory carries
the TMA-fed Q/K/V/$G$ views, Q/dS peer exchange, statistics, and instruction
descriptors.  The kernel is already at the practical register boundary:
an additional score buffer is ruled out by TMEM capacity.  The original
adaptive mixed-dP attempt also aliased its fixed dP scale word with the live
adaptive score scale page, causing repeatable tile-local corruption.  Giving
those values distinct lifetime-safe pages produces a stable 128-register,
zero-spill specialization.  Useful further pipelining still needs a different
ownership topology or an upstream/downstream fusion, not another local barrier
stage.

\subsection{Only dQ is a global reduction}

For an owner $o$, the kernel forms a partial
\begin{equation}
  dQ^{(o)}=\alpha\,dS^{(o)}K^{(o)},\qquad
  dQ=\sum_o dQ^{(o)}.
  \label{eq:dq-owner-reduction}
\end{equation}
The retained implementation publishes these partials with
\code{cp.reduce.async.bulk.tensor} addition.  It is not a scalar atomic loop.
At S8192/H8 the FP32 version moved 830.5 MB into the reduction path, 16.5
times the final 50.3 MB FP32 dQ tensor.  The retained BF16 reduction halves
that payload.  dK and dV stay in TMEM for their complete owner lifetime and
are stored once, so there is no dK/dV atomic mechanism to replace.
